\subsection{What this audit shows, and what it does not}

Two markers (grade, phenotype) survive the strongest test we apply -- zero-shot transfer to a
different institution and scanner with no retraining -- and, for grade, further survive
transfer to a completely different cohort against real pathologist ground truth (PRECISE,
$\rho=0.87$). PTEN loss is associated with its target and stable across TCGA-PRAD sites, but
its nested held-out increment over grade is small and its bootstrap interval includes zero --
a pattern that holds for CONCH and, independently, for Virchow, so it is not an artifact of one
encoder's embedding space. One (SPOP) is a clean, cross-encoder-confirmed null, consistent with
our own direct test and refutation of the two most likely alternative explanations (random-tile
dilution, resolution artifacts) for its disagreement with a prior single-institution report. AR
activity is associated in pooled internal CV but site-unstable, and its nested grade-adjusted
increment is likewise uncertain and likewise encoder-convergent. We regard this combination --
positive, negative, and unstable results, each supported by the same battery of checks -- as
more informative than a marker pool curated to show only positive findings would be.

What the audit does \emph{not} show is that qualification for one target implies clinical
utility for a different one. The LEOPARD pool comparison was prospectively frozen to test
exactly this: whether markers qualified against grade, phenotype, PTEN, or AR predict a real,
independent, downstream clinical outcome (biochemical recurrence) better than an unqualified
or null-inclusive pool. They do not -- every pool scored at or below chance. We consider this
a genuine, informative negative result rather than a failure to report: it directly
distinguishes ``this marker measures its stated target reliably'' from ``this marker is ready
to inform a clinical outcome model,'' a distinction that is easy to blur when validation
studies stop at the first (grade-, phenotype-, or molecular-status-level) demonstration.

\subsection{Reusable pitfalls}

Across this audit and its immediate predecessor experiments, we encountered the same handful
of failure modes repeatedly enough to consider them general properties of frozen pathology
foundation-model embeddings rather than one-off mistakes:

\begin{itemize}
\item \textbf{Tile-scale mismatch silently collapses transfer.} A probe trained at one
physical tile scale, applied to tiles at a substantially different scale, can fail completely
even when the underlying embedding space is not degenerate -- confirmed independently in two
unrelated transfers (an early NADT$\to$PANDA Virchow experiment, and this study's marker-7
Virchow transfer, where re-matching scale did \emph{not} recover performance, ruling scale
back \emph{out} as the explanation there and pointing instead to encoder-direction mismatch).
CONCH tolerated a comparable mismatch that broke Virchow, i.e.\ scale sensitivity itself is
encoder-dependent and not knowable in advance without testing.
\item \textbf{Different encoders can transfer in different directions even when both contain
real signal.} Marker 7 is the clearest example: CONCH transfers zero-shot from LEOPARD to
TCGA-PRAD; Virchow does not, despite an independently strong Virchow fit on \emph{each} cohort
alone and strong cross-encoder agreement on the LEOPARD in-cohort signal itself. A
single-encoder cross-cohort validation cannot distinguish ``the marker is real'' from ``this
particular encoder's fitted direction happens to transfer.''
\item \textbf{Qualification against one target does not certify transfer to a different
downstream outcome}, even when the outcome's signal is demonstrably present in the same
embedding space (Section~\ref{sec:leopard-results}). This is, to our knowledge, not
previously demonstrated directly in the literature we compare against.
\item \textbf{Hyperparameter choices can flip borderline results.} Standardizing a single
protocol across all 13 tested hypotheses changed one candidate (TCGA-PRAD H\&E$\to$ERG fusion
status) from nominally significant to null; we adopt the standardized number as authoritative
throughout and flag this as a reason to distrust single-protocol borderline claims in general.
\item \textbf{Label and outcome definitions are not interchangeable across cohorts even under
the same name.} TCGA-PRAD's GDC-harmonized clinical data has no populated
\texttt{days\_to\_recurrence} field for this cohort; a recurrence-style label reconstructed
from \texttt{follow\_ups.disease\_response} is a coarser, clinical-assessment-based construct
than LEOPARD's explicit PSA-based biochemical-recurrence definition, and marker 7's TCGA-PRAD
transfer should be read with that caveat rather than as a strict replication of the same
outcome.
\end{itemize}

\subsection{Relationship to prior confounder-audit and benchmarking work}

Dawood et al.\ perform a similarly motivated, larger-scale confounder audit but do not include
prostate cancer, do not cross-validate across independently trained encoders, and do not test
whether qualification changes a marker's value for a downstream clinical outcome -- the
comparison we make directly in Section~\ref{sec:leopard-results}. Gevaert/Bareja et al.\ and
Neidlinger et al.\ benchmark many encoders across many tasks, including several of the exact
targets we study (TCGA-PRAD ERG status, SICAPv2 grading), but their reported methodology does
not appear to enforce patient-disjoint evaluation, which our own early results show is
necessary to avoid a specific, previously undetected pseudo-replication failure mode. Our
contribution is narrower than either in absolute scope (prostate-specific, seven markers) but
adds three things neither provides together: mandatory patient-disjoint discipline throughout,
cross-encoder replication of the entire marker pool rather than a subset, and a direct
same-outcome test of whether qualification is consequential rather than merely descriptive.

\subsection{Limitations}

We did not acquire an independent molecular-validation cohort (new IHC or sequencing);
PTEN and AR-activity's molecular concordance rests on already-public TCGA-PRAD multi-assay
data, and ERG's biological mechanism (why residual signal predicts grade at all) remains an
open, unverified hypothesis. LEOPARD's public labels carry no clinical covariates, so an
incremental-value-over-clinical-baseline claim is untestable with this cohort; the pool
comparison in Section~\ref{sec:leopard-results} is accordingly scoped to outcome association,
not incremental clinical value. Marker 7 is a post-hoc discovery, not a protocol-frozen
candidate, and its TCGA-PRAD outcome definition differs in granularity from LEOPARD's; among
censored TCGA-PRAD patients, marker 7's risk score correlates with follow-up duration more
than the corresponding LEOPARD check did ($\rho=-0.28$, $p=3\times10^{-5}$ vs.\ $p=0.07$ on
LEOPARD), a residual concern we could not fully resolve given TCGA-PRAD's multi-institution
heterogeneity and coarser outcome definition, though the landmark analysis argues against a
pure observation-time artifact. Two further limitations emerged from expanding this audit in
response to external review (Section~\ref{sec:res-marker7-clinical},
\ref{sec:res-label-benchmark}): marker 7's added value over grade does not survive once a
fuller clinical covariate set (age, T-stage, PSA, surgical margin) is included, on a smaller
complete-case subset where margin status is itself a very strong predictor -- so ``grade-
independent'' should not be read as ``independent of standard clinical covariates generally'';
and re-scoring marker 7 directly against TCGA's standardized PFS/DFS endpoints, rather than our
reconstructed label, gives a visibly attenuated (though still above-chance) transfer, indicating
the original headline C-index is partly specific to our label's construction. What specific
morphological content the shared CONCH/Virchow
recurrence-predictive direction reflects remains unresolved: a preliminary, qualitative
tile-level review ($n=6$ patients, 18 tiles, no pathologist review) suggested loss of organized
benign stromal architecture as a plausible pattern, but one high-scoring tile showed an
atypical color cast raising the possibility of a staining or scanning artifact, and we present
this only as hypothesis-generating.

\subsection{Conclusion}

A statistically significant association between a pathology foundation-model embedding and a
clinical or molecular label is a starting point, not an endpoint. Applied with prospectively
frozen criteria to six initial candidates and transparently extended to one post-hoc marker,
this audit separates genuinely transportable signal (grade and phenotype) from provisional
or context-dependent findings (PTEN, the direct recurrence score, and AR) and a null result
(SPOP). In its most consequential test, qualification for one target is not evidence of
qualification for another. The frozen protocol and the pitfalls it
surfaced are offered as reusable tools for the next study asked to make the same distinction.
